\section*{图目录}
\begin{center}{\zihao{5} 图目录检索表}\end{center}
{\zihao{5}
\renewcommand{\arraystretch}{0.96}
\setlength{\LTpre}{0pt}
\setlength{\LTpost}{0pt}
\begin{longtable}{>{\centering\arraybackslash}p{1.8cm}>{\centering\arraybackslash}p{10.0cm}>{\centering\arraybackslash}p{1.5cm}}
\toprule
\textbf{图号} & \textbf{图名} & \textbf{页码} \\
\midrule
\endfirsthead
\toprule
\textbf{图号} & \textbf{图名} & \textbf{页码} \\
\midrule
\endhead
1-1 & 全文逻辑与计算模块关系 & 7 \\
2-1 & 三种候选车架方案原型照片 & 8 \\
2-2 & 三种候选方案综合评分对比 & 9 \\
2-3 & MIDAS 车身模型内力分布结果 & 11 \\
2-4 & 最终采用的空心双主梁勒勒车模型整体俯视示意 & 12 \\
3-1 & 单根主梁等效简支梁模型 & 13 \\
3-2 & 方形空心截面与等面积实心杆对比 & 14 \\
3-3 & 主梁壁厚对弯曲应力的影响 & 16 \\
3-4 & 布袋底面与多道横向小杆共同承托示意 & 17 \\
3-5 & 不同横向小杆布置的承托示意 & 18 \\
4-1 & 力的传递路径 & 19 \\
4-2 & 轮轴同轴度校正示意 & 21 \\
5-1 & 横杆端部胶接节点优化示意 & 22 \\
5-2 & 四螺丝牵引连接节点示意 & 23 \\
5-3 & 关键连接构造优化示意 & 24 \\
6-1 & 坡道牵引与布袋防滑示意 & 25 \\
6-2 & 弯道抗倾覆与重心投影示意 & 26 \\
6-3 & 动力放大系数与车轮反力关系 & 27 \\
6-4 & 单桥工况下车架抗扭示意 & 27 \\
6-5 & 典型工况风险分布图 & 28 \\
7-1 & 装载袋数对载质比的影响 & 31 \\
7-2 & 40 袋装货平面布置示意 & 31 \\
7-3 & 满载码放的重心控制示意 & 32 \\
8-1 & 现场制作时间甘特图 & 33 \\
8-2 & 制作完成后的校核流程图 & 34 \\
8-3 & 赛前核查卡片示意 & 36 \\
8-4 & 实物模型、加载试验和满载行驶照片 & 37 \\
9-1 & 作品设计、计算与试验闭环优化框架 & 39 \\
9-2 & 作品照片 & 40 \\
\bottomrule
\end{longtable}
}
